\section{Context-Free Grammars}
\subsection{Grammars and Languages}

\begin{definition}[Production/Rewrite Rule, Rewrite System]
    Given an alphabet \(\Omega\), a \term{production/rewrite rule} is an element \(\pqty{\alpha, \beta}\) of \(\Omega\kplus \times \Omega\kstar\), written as \(\alpha \prodrule \beta\), interpreted as ``replace \(\alpha\) with \(\beta\)''.

    A \term{rewrite system} \(\pqty{\Omega, P}\) is an alphabet \(\Omega\) and a finite set \(P\) of rewrite rules.
\end{definition}

Think of this as if \(w \in \Omega\kplus\) contains \(\alpha\) as a substring, so \(w = \gamma \alpha \delta\). Then we can \emph{choose} to replace \(\alpha\) with \(\beta\), changing \(w\) to \(\gamma \beta \delta\).

\begin{remark}
    We require \(\alpha \neq \emptyword\), but we allow \(\alpha \prodrule \emptyword\). This is known as an \term{\(\emptyword\)-production}.
\end{remark}

\begin{example*}
    Consider a group \(G = \angBrac{a, b}\) and \(\Omega = \Bqty{a, b, A, B}\). Then for any \(w \in \Omega\kplus\) containing \(bB\), applying a rewrite rule \(bB \prodrule \emptyword\) changes \(w\), but does not change the group element represented by \(w\) in \(G\).
    
    We can have the rewrite system
    \[
        P = \Bqty{aA \prodrule \emptyword, bB \prodrule \emptyword, Bb \prodrule \emptyword, Aa \prodrule \emptyword}.
    \]
    
    This is related to the idea of \term{group presentations} and \term{free groups}.
\end{example*}

\begin{notation}
    We often use capitals for non-terminals, lowercase for terminals, and lowercase Greek for words over \(\Omega\).
\end{notation}

But where is the language?

\begin{definition}[Variable-Based Grammar, Derivation]
    A \term{(variable-based) grammar} \(G = \pqty{\Sigma, N, P, S}\) is a rewrite system where:
    \begin{itemize}
        \item \(\Omega = \Sigma \sqcup N\) is divided into \term{terminal letters} \(\Sigma\) and \term{non-terminal letters/variables} \(N\);
        \item Any production rule \(\pqty{\alpha \prodrule \beta} \in P\) only has non-terminals on the left-hand side, i.e. \(\alpha \in N\kplus\), and \(P \subseteq N\kplus \times \Omega\kstar\);
        \item \(S \in N\) is a single \term{start symbol}.
    \end{itemize}

    For a grammar \(G = \pqty{\Sigma, N, P, S}\) and \(\alpha, \beta \in \Omega\kstar\), we say
    \begin{enumerate}[label=(\roman*)]
        \item \term{\(\beta\) is derivable from \(\alpha\) in one step} if \(\alpha = \alpha_1 \gamma \alpha_2\), \(\beta = \alpha_1 \delta \alpha_2\), and \(\pqty{\gamma, \delta} \in P\), written as \(\alpha \prodrule[1][G] \beta\).
        \item \term{A derivation of length \(n\) of \(\beta\) from \(\alpha\)}, written \(\alpha \prodrule[n][G] \beta\), is inductively defined as there exists some \(\gamma\) with \(\alpha \prodrule[n-1][G] \gamma\) and \(\gamma \prodrule[1][G] \beta\). Where \(n = 0\), we have \(\alpha \prodrule[0][G] \alpha\).
        \item \term{\(\beta\) is derivable from \(\alpha\)}, denotes as \(\alpha \prodrule[\ast][G] \beta\) means that there is some \(n \in \NN\), such that \(\alpha \prodrule[n][G] \beta\).
        \item A word \(w \in \Omega\kstar\) derivable for start symbol \(S\) is in \term{sequential form}.
        \item A word in sequential form is a \term{sentence}, if \(w \in \Sigma\kstar\) consists of only terminals.
    \end{enumerate}
\end{definition}

\begin{definition}[Language Generated by a Grammar]
    The \term{language generated by a grammar \(G\)} is the set
    \[
        \langacc\pqty{G} = \set{w \in \Sigma\kstar}{S \prodrule[\ast][G] w}.
    \]

    This is the set of all sentences derivable in \(G\) from the start symbol \(S\).
\end{definition}

\begin{remark}
    In general, a word \(w \in \langacc\pqty{G}\) could have many derivations from \(S\), of same or even different lengths. Also, we can have grammars \(G_1 \neq G_2\) with \(\langacc\pqty{G_1} = \langacc\pqty{G_2}\).

    For example, if we have a non-terminal character \(X \neq S\) not appearing in the right-hand side of any production rule \(\alpha \prodrule \beta \in P\). Then \(X\) never appears in any derivation from \(S\), so we can add \(X \prodrule \emptyword\) to \(P\) which is never used.
\end{remark}

The \term{Chomsky hierarchy} classifies grammars and languages as below.

\begin{definition}[Type of Production Rule, Grammar and Language]
    A production rule \(\pqty{\alpha, \beta} \in \Omega\kplus \times \Omega\kstar\) is of:
    \begin{itemize}
        \item \term{Type 0} if \(\alpha \in N\kplus\);
        \item \term{Type 1 (Noncontracting)} if \(\alpha \in N\kplus\) and \(\abs{\alpha} \leq \abs{\beta}\);
        \item \term{Type 2} if \(\alpha = A \in N\) and \(\abs{\beta} \geq 1\);
        \item \term{Type 3} if \(\alpha = A \in N\) and \(\beta = a\) or \(\beta = aB\) for \(B \in N\) and \(a \in \Sigma\).
    \end{itemize}

    A \term{grammar} is of \term{type \(i\)} if all its production rules are of type \(i\). A \term{language} is of \term{type \(i\)} if it is generated by some grammar of type \(i\).
\end{definition}

\begin{remark}
    A grammar of type 0 is any variable-based grammar.
\end{remark}

Suppose \(L\) is of type 1 over \(\Sigma\) and we are given \(w \in L\). Then there exists a derivation in type 1 grammar \(G\), such that
\[
    \beta_0 = S \prodrule[1][G] \beta_1 \prodrule[1][G] \beta_2 \prodrule[1][G] \cdots \prodrule[1][G] \beta_n = w,
\]
and we take the shortest such derivation, with length \(m\).

But \(\abs{\beta_i} \leq \abs{\beta_{i + 1}}\), and if \(l = \abs{\beta_i} = \cdots = \abs{\beta_{i + j}}\) for \(j \geq \abs{\Omega}^l\), then two of these \(\beta\)s are the same by the pigeonhole principle, so \(m\) must not be minimal.

The, we have
\[
    m \leq \sum_{l = 1}^{\abs{w}} \abs{\Omega}^l = M
\]
so we can tell whether or not \(w \in L = \langacc\pqty{G}\) systematically by creating all derivations of length at most \(M\) from \(S\).

If \(\emptyword \in \langacc\pqty{G}\), then we need some \(\emptyword\)-production \(\alpha \prodrule \emptyword\), so languages of type 1 (and hence, of 2 and 3) do \emph{not} contain the empty word.

\subsection{Context-Free Languages}

\begin{definition}[Context-Free Grammar and Language]
    A \term{context-free grammar} (CFG) \(G = \pqty{N, \Sigma, P, S}\) is where every production rule is of the fom \(\pqty{A, \alpha} \in N \times \Omega\kstar\), i.e., a type 2 grammar but also allowing \(\alpha = \emptyword\). A \term{language} \(L\) is \term{context-free} (CFL) if it is generated by some context-free grammar.
\end{definition}

\begin{examples}[Context-Free Languages]
    \begin{enumerate}[label=(\roman*)]
        \item If a language \(L\) is of type 2, then it is context-free.
        \item Let \(N = \Bqty{S}\) and \(\Sigma = \Bqty{a, b}\), with production rules \(S \prodrule aSb\) and \(S \prodrule \emptyword\) which we denote as \(S \prodrule aSm \mid \emptyword\). Then we obtain a context-free language \(\emptyword, ab, aabb, \ldots, a^n b^n, \ldots\), which is \emph{not} regular!
        \item All regular languages are context-free.
        \item Consider the production rule
        \[
            S \prodrule SaSbS \mid SbSaS \mid \emptyword.
        \]

        We get words with the same number of \(a\)s and \(b\)s.
        \item For \(w = \sigma_1 \cdots \sigma_n \in \Sigma\kstar\), the \term{reverse} of \(w\) is \(R\pqty{w} = \sigma_n \cdots \sigma_1 \in \Sigma\kstar\). The set of palindromes
        \[
            X = \set{w \in \Bqty{a, b}\kstar}{R\pqty{w} = w}
        \]
        is a context-free language. It is generated by the production rules
        \[
            S \prodrule aSa \mid bSb \mid a \mid b \mid \emptyword.
        \]

        The first two production rules give symmetry ``from the middle'', and the next to finish it off if the final length is odd, and the empty word finishes it off if the final length is even.
    \end{enumerate}
\end{examples}

\begin{definition}[Language of Balanced Parentheses]
    Let \(\Sigma = \Bqty{[, ]}\). For \(w \in \Sigma\kstar\), let \(L\pqty{w}\) be the number of \([\)s in \(w\) and \(R\pqty{w}\) be the number of \(]\)s in \(w\). We say \(w\) is \term{balanced} if \(L\pqty{w} = R\pqty{w}\) and for all \(x \in \Sigma\kstar\) with \(w = xy\), we have \(L\pqty{x} \geq R\pqty{x}\).
\end{definition}

\begin{proposition}[Grammar of Balanced Parentheses]
    The context-free grammar \(G\) with production rule
    \[
        S \prodrule[][] [S] \mid SS \mid \emptyword
    \]
    has \(\langacc\pqty{G}\) as the language of all balanced words in \(\Sigma\kstar\). This is \emph{not} regular.
\end{proposition}